\section{Responses to Anticipated Critiques}
\label{sec:responses}

Because the X--$\Theta$ framework departs from standard path-independence assumptions
commonly adopted in Bell analyses, it is natural to anticipate several conceptual objections.
We address the most critical concerns explicitly.

\subsection{No-Signaling and the Absence of Hidden Communication}

\paragraph{Concern.}
The introduction of a connection $\phi(a,b)$ over detector settings space may appear
to enable superluminal signaling: a change in Alice’s setting could, in principle,
alter the effective phase experienced at Bob’s site.

\paragraph{Resolution.}
X--$\Theta$ enforces no-signaling at the level of operational marginals.
Local outcomes are generated by symmetric thresholding of an internal phase variable
$\Theta_{\mathrm{eff}} = \Theta + \phi(a,b)$ with $\Theta$ sampled uniformly from $[0,2\pi)$.
As a result,
\begin{equation}
P(x|a,b) = \int d\Theta\, P(x|a,\Theta_{\mathrm{eff}}) = P(x|a),
\end{equation}
and analogously for Bob.
The connection modifies only \emph{relative phases} appearing in joint correlators,
not local outcome statistics.

Operationally, the connection functions as a gauge field over configuration space.
As in the Aharonov--Bohm effect, gauge-invariant loop observables may change while
local measurements remain unaffected.
Any deviation from marginal flatness would constitute a failure mode of the framework
and is explicitly checked in simulation and proposed experiments.

\subsection{Parameter Dependence and the Allegation of Fine-Tuning}

\paragraph{Concern.}
The appearance of Tsirelson-saturating correlations may be attributed to tuning of
model parameters such as detector saturation or visibility.

\paragraph{Resolution.}
The parameters appearing in X--$\Theta$ represent physical features of real detectors,
not arbitrary fitting knobs.
Nonlinear thresholding is unavoidable in discrete-outcome measurements and acts
as a harmonic filter on the underlying phase variable.
Under mild regularity assumptions, the lowest harmonic dominates:
\begin{equation}
E(a,b) \propto -\cos(a-b+\beta),
\end{equation}
independently of microscopic details.

The Tsirelson bound arises here as a geometric ceiling: it corresponds to the maximal
holonomy attainable on a four-vertex plaquette under a $U(1)$ connection while preserving
no-signaling.
Attempts to exceed $2\sqrt{2}$ require either higher-curvature transport or explicit
dependence of marginals on distant settings, neither of which is permitted by the axioms.

\subsection{Contextuality Without Nonlocal Knowledge}

\paragraph{Concern.}
Because $\Theta_{\mathrm{eff}}$ depends on both settings $(a,b)$, the model may be
interpreted as a contextual hidden-variable theory in which particles ``know'' distant
choices.

\paragraph{Resolution.}
The phase $\Theta$ is not carried by the particle as a message.
Instead, it couples locally to a geometric structure defined over the
\emph{settings manifold}.
Contextuality enters through path-dependent holonomy, not through instantaneous
setting awareness.

In Version-0B, the accumulated phase depends on the orientation of traversal through
configuration space:
\begin{equation}
\Omega_{\gamma} = \sum_{t \in \gamma} \delta\, \sigma_t,
\end{equation}
where $\sigma_t$ encodes local orientation.
The particle interacts only with a local internal degree of freedom that has been
geometrically transported by prior apparatus motion.
No assumption of foreknowledge or superluminal update is required.

\subsection{Why Path-Dependent Effects Have Escaped Detection}

\paragraph{Concern.}
If a CW/CCW asymmetry $\Delta S$ exists, why has it not appeared in decades of Bell tests?

\paragraph{Resolution.}
Standard Bell experiments are explicitly designed to eliminate memory effects by
randomizing settings at each trial.
This procedure symmetrizes over orientation in configuration space, averaging
clockwise and counter-clockwise contributions:
\begin{equation}
\langle \Delta S \rangle_{\mathrm{randomized}} \rightarrow 0.
\end{equation}

X--$\Theta$ predicts that $\Delta S$ becomes observable only under controlled,
non-random traversal of the CHSH loop.
This renders the effect invisible to conventional analyses but accessible through
re-analysis of time-ordered data or through dedicated cycling protocols.

\paragraph{Falsifiability.}
If experiments conducted under controlled CW and CCW cycling yield $\Delta S$
consistent with zero while maintaining strict no-signaling marginals,
the framework is falsified.

\subsection{Summary}

The X--$\Theta$ framework preserves all operational constraints of Bell experiments
while relaxing a single assumption: path-independence in detector settings space.
Its predictions are geometric, bounded, and experimentally testable.
No superluminal signaling, parameter fitting, or unphysical contextual knowledge
is required.
